\documentclass[12pt]{article}
\usepackage[utf8]{inputenc}
\usepackage[T1]{fontenc}
\usepackage{hyperref}
\title{Tareng Language Resources for K-12 Curriculum Development}
\author{Compiled by Claude Sonnet 5 (claude-sonnet-5) for Andr\'as Kornai}
\date{}
\begin{document}
\maketitle

\section{Introduction}
Tareng (ISO 639-3 tgr; Glottocode tare1239) is spoken in Laos. Its classification is disputed. Ethnologue and Joshua Project treat it as Katuic. Glottolog says that placement is probably wrong, and its entry describes the language as Bahnaric and says the Tareng and Talieng entries may reflect a Katuic/Bahnaric confusion. Other names are Tariang and Laoseng. Joshua Project lists about 12,000 speakers, mostly in Phongsali province, and Global Recordings gives 11,000. The ``Bookkeeping'' family label in the task header appears to be a Glottolog artefact. The search results suggest the language is still spoken, so the heritage framing may not apply.

\section{Primary Educational Resources}
No primers, textbooks or school programmes were found.

\section{Linguistic and Reference Materials}
\begin{itemize}
\item Glottolog entry for Tareng, with the retirement and classification note: \url{https://glottolog.org/resource/languoid/iso/tgr}
\end{itemize}
No dictionary or grammar was found.

\section{Religious and Community Resources}
\begin{itemize}
\item Joshua Project profile of the Laoseng (Tareng), which reports that Bible translation has started: \url{https://joshuaproject.net/people_groups/15272}
\end{itemize}

\section{Digital and Audio Resources}
\begin{itemize}
\item Global Recordings Network page for tgr (audio): \url{https://globalrecordings.net/en/language/tgr}
\end{itemize}
I did not verify the content or the variety of these recordings.

\section{Curriculum Development Context}
Almost nothing exists. A first step would be to settle the identity and classification of the language, and then to document it with the community.
\end{document}